
\section{边界附近的椭圆估计及其应用}
\label{sec:elliptic-estimates-boundary}

\hzsync{theorem}{10}%
\hzsync{equation}{42}%

把由
\[
B_{\varepsilon}^{+}\ =\ \bigl\{\,z\in\C\ \bigm|\ |z|<\varepsilon,\ \operatorname{Re}(z)\geq0\,\bigr\}
\]
给出的\trHalfBall{}记作 $B_{\varepsilon}^{+}$，其\trBoundary{}为%
 $\partial B_{\varepsilon}^{+}=(-\varepsilon,\varepsilon)i$。

对 $k\in\mathbb{N}$ 与 $1<p<\infty$，把通常的\trSobolev{}\trNorm{}记作%
 $\|\cdot\|_{k,p,\Omega}$，即
\[
\|u\|_{k,p,\Omega}^{p}\ =\ \sum_{|\alpha|\leq k}\ \bigl\|D^{\alpha}u\bigr\|_{L^{p}(\Omega)}^{p}\,.
\]
对于 $\bar{\partial}=\frac{\partial}{\partial s}+i\frac{\partial}{\partial t}$，%
经典的估计（例如可在\trVekua{} [215] 中找到）即是下述

\begin{theorem}\label{thm:half-ball-elliptic-estimate}
给定 $\varepsilon>0$、$k\in\mathbb{N}$ 与 $1<p<\infty$，存在常数%
 $c=c(k,p,\varepsilon)>0$，使得
\[
\bigl\|\bar{\partial}u\bigr\|_{k,p,B_{\varepsilon}^{+}}\ \geq\ c\,\|u\|_{k+1,p,B_{\varepsilon}^{+}}
\]
对一切光滑映射 $u:B_{\varepsilon}^{+}\to\C^{n}$ 成立，此处 $u$ 在%
 $B_{\varepsilon}^{+}$ 内有\trCompactSupport{}并满足\trBoundaryCondition
\[
u\bigl(\partial B_{\varepsilon}^{+}\bigr)\ \subset\ \R^{n}\,.
\]
\end{theorem}

利用这一估计以及标准\trEllipticRegularity{}理论中常用的\trDifferenceQuotient{}技巧，%
我们得到下述正则性结果。

设 $J$ 是定义在 $\C^{n}$ 上（或至少在 $0$ 附近）的一个\trSmooth%
\trAlmostComplexStructure，满足 $J(0)=i$。设
\[
u\ :\ B_{\varepsilon}^{+}\ \longrightarrow\ \C^{n}
\]
属于 $W^{1,p}$ 类，此处 $p>2$，并且满足
\begin{align*}
u(0)\ &=\ 0\\
u\Bigl((-\varepsilon,\varepsilon)i\Bigr)\ &\subset\ \R^{n}\\
\frac{\partial u}{\partial s}\ +\ J(u)\frac{\partial u}{\partial t}\ &=\ 0\qquad\text{于 }\mathring{B}_{\varepsilon}^{+}
\end{align*}
此处 $\mathring{B}_{\varepsilon}^{+}=B_{\varepsilon}^{+}\setminus(-\varepsilon,\varepsilon)i$。我们有下述

\begin{theorem}\label{thm:boundary-smoothness}
在上述假设之下，存在 $\delta\in(0,\varepsilon)$，使得 $u|B_{\delta}^{+}$ 是%
\trSmooth{}的。
\end{theorem}

我们的下一个目标是对第 6 章中提出的一个正则性问题给出一个应用。设
\[
w\ :\ [-\varepsilon,\varepsilon]\times[-\varepsilon,\varepsilon]\ \longrightarrow\ \C^{n}
\]
属于 $W^{1,p}$ 类，此处 $2<p<\infty$。设 $w(0)=0$，而 $J$ 是定义在零点的一个%
充分大的\trNeighborhood{}上的\trAlmostComplexStructure，使得 $J(0)=i$。再设%
 $H:[-\varepsilon,\varepsilon]\times U\to\R$ 是一个\trSmooth{}映射\footnote{原书印作「$H:[-\varepsilon,\varepsilon]\times U\to\C^{n}$」：然此处 $H'$ 乃 $H$ 关于第二个变量取的梯度（梯度只对实值函数有定义），且 $H'(t,w)$ 以 $\C^{n}$ 中向量之身份进入方程 \eqref{eq:a42-perturbed-cr}，故 $H$ 必须实值，$\C^{n}$ 系 $\R$ 之误，此处径改。——译者注}，其中%
 $U$ 是 $\C^{n}$ 中 $0$ 的一个充分大的\trNeighborhood。此外，令 $H'(0,0)=0$，%
这里 $H'$ 是关于第二个变量、相对于 $U$ 上某个\trRiemann{}度量取的\trGradient。

我们考虑\trPartialDifferentialEquation
\begin{equation}\label{eq:a42-perturbed-cr}
\frac{\partial w}{\partial s}\ +\ J(w)\frac{\partial w}{\partial t}\ +\ \chi_{[0,\infty)}(s)H'(t,w)\ =\ 0
\end{equation}
于 $(-\varepsilon,\varepsilon)\times(-\varepsilon,\varepsilon)$ 上。此处%
 $\chi_{[0,\infty)}$ 是 $[0,\infty)$ 的示性函数。令人称奇的事实是：对某个%
充分小的 $\delta\in(0,\varepsilon)$，
\[
w|\Bigl([-\delta,0]\times[-\delta,\delta]\Bigr)
\]
与
\[
w|\Bigl([0,\delta]\times[-\delta,\delta]\Bigr)
\]
都是\trSmooth{}的。特别地，
\[
w|\Bigl(\{0\}\times[-\delta,\delta]\Bigr)\ \longrightarrow\ \C^{n}
\]
是\trSmooth{}的。注意 $w|\bigl([-\delta,\delta]\times[-\delta,\delta]\bigr)$ %
一般并不是\trSmooth{}的，因为 $s$ 方向的导数未必相匹配。上述光滑性断言的证明将通过%
一系列\trTransformation{}化归为\cref{thm:boundary-smoothness}。

为使论证简单，我们设所有数据 $H$ 与 $J$ 都定义在整个 $\C^{n}$ 上。在%
 $[0,\varepsilon]\times[-\varepsilon,\varepsilon]\times\C^{n}\times\C^{n}$ 上%
如下定义\trAlmostComplexStructure{} $\hat{J}$：
\[
\hat{J}(s,t,u,v)(a,b,h,k)\ =\ \Bigl(\,-b,\ a,\ J(u)h\ +\ bH'(t,u)\ +\ aJ(u)H'(t,u),\ -J(v)k\Bigr)\,,
\]
$\hat{J}$ 的这一定义受\trGromov{}技巧以及一种人们不妨称之为「逆反射原理」%
（inverse reflection principle）的东西的启发。我们将看到，%
\eqref{eq:a42-perturbed-cr} 的一个解 $w$ 给出结构 $\hat{J}$ 之下%
 $[0,\varepsilon)\times[-\varepsilon,\varepsilon)$ 上的一个\trHolomorphic{}%
映射 $\alpha$。

我们留给读者去验证 $\hat{J}^{2}=-\mathrm{Id}$。令
\[
j\ =\ \hat{J}(0,0,0,0)\,.
\]
我们观察到
\[
j(a,b,h,k)\ =\ \Bigl(\,-b,\ a,\ ih\ +\ bH'(0,0)\ +\ aiH'(0,0),\ -ik\Bigr)\,.
\]
如下定义 $\R\times\R\times\C^{n}\times\C^{n}$ 的一个实\trLinear{}子空间：
\[
T\ =\ \bigl\{\,(0,b,h,h)\ \bigm|\ b\in\R,\ h\in\C^{n}\,\bigr\}\,.
\]
显然 $\dim_{\R}(T)=2n+1=:N$。我们注意到
\begin{equation}\label{eq:a43-total-real-t}
T\ \cap\ jT\ =\ \{0\}\,.
\end{equation}
的确，若
\[
(0,b_{1},h_{1},h_{1})\ =\ j(0,b_{2},h_{2},h_{2})\ =\ \Bigl(\,-b_{2},\ 0,\ ih_{2}+b_{2}H'(0,0),\ -ih_{2}\Bigr),
\]
则得 $b_{1}=b_{2}=0$，这意味着
\begin{align*}
ih_{2}\ &=\ h_{1}\\
-ih_{2}\ &=\ h_{1}\,,
\end{align*}
从而 $h_{1}=h_{2}=0$。如下定义映射%
 $\alpha:[0,\varepsilon]\times[-\varepsilon,\varepsilon]\to\R\times\R\times\C^{n}\times\C^{n}$\footnote{原书此应用部分通篇把未知函数印作「$u$」，系「$w$」之误——本部分所设未知量自始至终为 $w$（\eqref{eq:a42-perturbed-cr} 之解），而 $u$ 乃\cref{thm:boundary-smoothness} 陈述中之映射名——以下各式均径改作 $w$，不另注。——译者注}：
\[
\alpha(s,t)\ =\ \Bigl(\,s,\ t,\ w(s,t),\ w(-s,t)\Bigr)\,.
\]
于是
\[
\alpha(0,t)\ =\ \Bigl(\,0,\ t,\ w(0,t),\ w(0,t)\Bigr)\in T
\]
对 $t\in[-\varepsilon,\varepsilon]$ 成立。此外，由 $w$ 按假设具有 $W^{1,p}$ 正则性%
以及其余坐标分量的光滑性可知，$\alpha$ 属于%
 $W^{1,p}\bigl([0,\varepsilon]\times[-\varepsilon,\varepsilon],\,\R\times\R\times\C^{n}\times\C^{n}\bigr)$。%
在 $(0,\varepsilon)\times(-\varepsilon,\varepsilon)$ %
上我们算得
\begin{align*}
&\alpha_{s}(s,t)\ +\ \hat{J}\bigl(\alpha(s,t)\bigr)\bigl(\alpha_{t}(s,t)\bigr)\\
&\qquad=\ \Bigl(\,1,\ 0,\ w_{s}(s,t),\ -w_{s}(-s,t)\Bigr)\\
&\qquad\quad+\ \Bigl(\,-1,\ 0,\ J(w(s,t))w_{t}(s,t)\ +\ H'(t,w(s,t)),\ -J(w(-s,t))w_{t}(-s,t)\Bigr)\\
&\qquad=\ (0,0,0,0)\,.
\end{align*}
鉴于 \eqref{eq:a43-total-real-t}，集合 $T$ 是 $(\R\times\R\times\C^{n}\times\C^{n},j)$ %
的一个\trTotallyReal{}子空间。找一个复\trLinear{}映射
\[
\Phi\ :\ \bigl(\R\times\R\times\C^{n}\times\C^{n},j\bigr)\ \longrightarrow\ \bigl(\C^{N},i\bigr)
\]
把 $T$ 映满 $\R^{N}$，这是一道简单的练习。回忆 $N=2n+1$。定义 $\beta=\Phi\circ\alpha$ %
与 $\C^{N}$ 上的一个\trAlmostComplexStructure{} $\tilde{J}$ 如下：
\[
\tilde{J}(A)\ =\ \Phi\hat{J}\Bigl(\Phi^{-1}(A)\Bigr)\Phi^{-1}\,.
\]
我们求得
\begin{align*}
\tilde{J}(0)\ &=\ \Phi j\Phi^{-1}\\
&=\ i\Phi\Phi^{-1}\\
&=\ i\,.
\end{align*}
此外，在 $(0,\varepsilon)\times(-\varepsilon,\varepsilon)$ 上，
\begin{align*}
\beta_{s}\ +\ \tilde{J}(\beta)\beta_{t}\ &=\ \Phi\alpha_{s}\ +\ \Phi\hat{J}\bigl(\Phi^{-1}\Phi\alpha\bigr)\Phi^{-1}\Phi\alpha_{t}\\
&=\ \Phi\bigl[\alpha_{s}\ +\ \hat{J}(\alpha)\alpha_{t}\bigr]\\
&=\ 0\,.
\end{align*}
最后，我们观察到
\[
\beta(0,t)\in\R^{N}
\]
对 $t\in(-\varepsilon,\varepsilon)$ 成立。应用\cref{thm:boundary-smoothness} %
我们推出：由于 $\beta\in W^{1,p}\bigl([0,\varepsilon]\times[-\varepsilon,\varepsilon],\C^{N}\bigr)$ %
且 $\beta\bigl(i(-\varepsilon,\varepsilon)\bigr)\subset\R^{N}$，故对某个%
 $\delta\in(0,\varepsilon)$，$\beta$ 在 $[0,\delta]\times[-\delta,\delta]$ 上是%
\trSmooth{}的。于是
\[
[0,\delta]\times[-\delta,\delta]\ \longrightarrow\ \C^{n}\ :\ (s,t)\ \longmapsto\ w(s,t)
\]
与
\[
[0,\delta]\times[-\delta,\delta]\ \longrightarrow\ \C^{n}\ :\ (s,t)\ \longmapsto\ w(-s,t)
\]
都是\trSmooth{}的，此即所欲证。

我们以一个在第 6 章中用到的观察来作结。回忆第 6 章中 $S_{T}$ 的定义，并设%
 $w:S_{T}\to M$ 求解\footnote{原书此式印作「$H'(t,w)\,dt-J(u)H'(t,w)\,ds$」：其一 $J(u)$ 之 $u$ 系 $w$ 之误；其二 $ds$／$dt$ 之位置与第 6 章同款方程（$H'$ 配 $ds$、$JH'$ 配 $dt$）相反，且唯如此方与下文「\eqref{eq:a44-st-loop} 在这些坐标下像 \eqref{eq:a42-perturbed-cr}」之断言相容——均径改。——译者注}
\begin{equation}\label{eq:a44-st-loop}
0\ =\ \bar{\partial}w\ +\ \chi_{[-T,T]\times S^{1}}\Bigl(H'(t,w)\,ds\ -\ J(w)H'(t,w)\,dt\Bigr)\,.
\end{equation}
在 $Z_{T}\subset S_{T}$ 中一点 $(T,t_{0})$（相应地，$(-T,t_{0})$）附近\footnote{原书四处坐标点均印作「$(T_{0},t_{0})$」「$(-T_{0},t_{0})$」，末句并作「由于 $T_{0}$ 是任取的」。然 \eqref{eq:a44-st-loop} 中截断项 $\chi_{[-T,T]\times S^{1}}$ 恰在 $s=\pm T$ 的粘合圆处不连续，欲化归为 \eqref{eq:a42-perturbed-cr} 须在 $(\pm T,t_{0})$ 附近取坐标；而 $T$ 既已固定，使两条闭路整圆光滑者乃边界点参数 $t_{0}$ 的任意性。故前四处 $T_{0}$ 系 $T$ 之误、末句 $T_{0}$ 系 $t_{0}$ 之误，均径改。——译者注}%
取适当的复坐标，并在 $M$ 中于 $w(T,t_{0})$（相应地，$w(-T,t_{0})$）%
附近取适当的辛坐标，我们看到 \eqref{eq:a44-st-loop} 在这些坐标下看上去就像%
 \eqref{eq:a42-perturbed-cr}。特别地，由于 $t_{0}$ 是任取的，我们推出闭路%
 $t\mapsto w(-T,t)$ 与 $t\mapsto w(T,t)$ 都是\trSmooth{}的。
